\documentclass[11pt]{article}
\usepackage[utf8]{inputenc}
\usepackage[margin=1in]{geometry}
\usepackage{amsmath}
\usepackage{graphicx}
\usepackage{booktabs}
\usepackage{hyperref}
\usepackage[round]{natbib}
\usepackage{lineno}
\usepackage{xcolor}
\usepackage{multirow}
\usepackage{float}

\linenumbers

\title{Female-Specific Cortical Circuit Hyperexcitability Driven by mGluR5--ER$\alpha$ Signaling in a \textit{Pten} Deletion Model of Autism}

\author{
Author A$^{1,*}$, Author B$^{1}$, Author C$^{1}$, Author D$^{2}$, Author E$^{1}$ \\
\\
$^{1}$Department of Neuroscience, University of Texas Southwestern Medical Center \\
$^{2}$Department of Pharmacology, University of Texas Southwestern Medical Center \\
$^{*}$Corresponding author: authora@utsouthwestern.edu
}

\date{}

\begin{document}
\maketitle

\begin{abstract}
Autism spectrum disorder (ASD) is diagnosed four times more frequently in males, yet females with ASD often present with more severe sensory processing deficits and epilepsy. Here we identify a female-specific cortical circuit dysfunction caused by deletion of the ASD-risk gene \textit{Pten} in neocortical pyramidal neurons. In acute thalamocortical slices from NSEPten knockout (KO) mice, spontaneous UP state duration was increased by 56\% selectively in females (layer 4 somatosensory cortex; sex $\times$ genotype interaction: $F(1,73) = 5.56$, $p < 0.05$ in L2/3; $F(1,71) = 7.33$, $p < 0.01$ in L5). The mGluR5 negative allosteric modulator MTEP (3~$\mu$M) rescued prolonged UP states exclusively in female Cre(+) mice, and mGluR5 protein was selectively elevated in female L5-enriched cortical lysates ($p < 0.01$, unpaired $t$-test). Two selective ER$\alpha$ antagonists (MPP and GNE, 1~$\mu$M each) rescued female hyperexcitability, and genetic reduction of ER$\alpha$ (Esr1$^{fl/+}$) in Pten KO neurons normalized UP state duration ($p < 0.01$), de novo protein synthesis ($p < 0.001$), and soma size ($p < 0.05$) only in females. Co-immunoprecipitation revealed elevated mGluR5--ER$\alpha$ complexes in female cortex. Female NSEPten KO mice exhibited deficits in auditory temporal processing (gap-ASSR), mGluR5-dependent seizures, and behavioral abnormalities. These findings were replicated in the clinically relevant germline Pten$^{+/-}$ model. We identify a sex-specific mGluR5--ER$\alpha$--ERK--protein synthesis signaling pathway that drives cortical hyperexcitability in females, providing a molecular mechanism for sex-differential ASD phenotypes.
\end{abstract}

\section{Introduction}

Autism spectrum disorder (ASD) affects approximately 1 in 54 children, with a consistently reported 4:1 male-to-female prevalence ratio \citep{autism2018prevalence}. However, accumulating evidence indicates that ASD in females may be underdiagnosed due to sex-differential symptom presentation, and when diagnosed, females often exhibit more severe comorbidities including sensory processing abnormalities and epilepsy \citep{werling2013sex}. The biological basis for these sex differences remains poorly understood.

Mutations in the phosphatase and tensin homolog (\textit{PTEN}) gene cause a spectrum of disorders collectively termed PTEN Hamartoma Tumor Syndromes, which include macrocephaly, intellectual disability, ASD, and epilepsy \citep{butler2005subset}. Mouse models with conditional or germline \textit{Pten} deletion recapitulate core features of these disorders, including social behavior deficits and altered brain connectivity \citep{kwon2006pten}. Notably, germline Pten heterozygous (Pten$^{+/-}$) mice exhibit female-specific social behavior deficits \citep{clipperton2012pten, page2009haploinsufficiency}, yet the underlying neural circuit dysfunction has not been explored.

In breast cancer, PTEN loss-of-function intersects with estrogen receptor alpha (ER$\alpha$) signaling to drive tumor growth via the PI3K/Akt/mTOR pathway \citep{lee2003pten}. Whether similar sex-dependent molecular interactions operate in the brain to produce sex-specific neural circuit phenotypes has not been tested. Meanwhile, the metabotropic glutamate receptor mGluR5 has emerged as a critical node in ASD-related circuit dysfunction: mGluR5 hyperactivity drives cortical hyperexcitability in Fragile X syndrome models \citep{huber2002altered, hays2011altered, gibson2008imbalance}, and genetic correction of mGluR5 signaling rescues multiple Fragile X phenotypes \citep{dolen2007correction}.

Here we use the neocortical UP state---a spontaneous persistent activity state driven by recurrent excitatory-inhibitory circuits \citep{shu2005properties}---as a sensitive readout of circuit excitability in thalamocortical slices from NSEPten KO and Pten$^{+/-}$ mice. We discover a female-specific hyperexcitability phenotype and trace it to an mGluR5--ER$\alpha$--ERK--protein synthesis signaling cascade, providing the first molecular mechanism for sex-specific effects of an ASD-risk gene on neural circuit function.

\section{Results}

\subsection{Female-Specific UP State Prolongation in NSEPten KO Mice}

We performed extracellular multiunit recordings of spontaneous UP states in acute thalamocortical slices from P18--25 NSEPten KO mice and Cre(-) littermate controls (Table~\ref{tab:upstate_duration}).

\begin{table}[H]
\centering
\caption{UP state duration across cortical layers in NSEPten KO mice. Values are mean $\pm$ SEM. Statistical analysis by two-way ANOVA (sex $\times$ genotype) with Sidak post-hoc tests. Percent change calculated relative to same-sex Cre(-) controls.}
\label{tab:upstate_duration}
\begin{tabular}{llccccc}
\toprule
\textbf{Layer} & \textbf{Sex} & \textbf{Cre(-) (ms)} & \textbf{Cre(+) (ms)} & \textbf{Change (\%)} & \textbf{$F$(sex$\times$geno)} & \textbf{$p$-value} \\
\midrule
\multirow{2}{*}{L4} & Female & $412 \pm 28$ & $643 \pm 38$ & $+56.1$ & \multirow{2}{*}{$4.82$} & \multirow{2}{*}{$< 0.05$} \\
 & Male & $398 \pm 31$ & $421 \pm 34$ & $+5.8$ & & \\
\midrule
\multirow{2}{*}{L2/3} & Female & $378 \pm 25$ & $562 \pm 36$ & $+48.7$ & \multirow{2}{*}{$5.56$} & \multirow{2}{*}{$< 0.05$} \\
 & Male & $385 \pm 29$ & $402 \pm 32$ & $+4.4$ & & \\
\midrule
\multirow{2}{*}{L5} & Female & $445 \pm 32$ & $691 \pm 42$ & $+55.3$ & \multirow{2}{*}{$7.33$} & \multirow{2}{*}{$< 0.01$} \\
 & Male & $438 \pm 30$ & $462 \pm 35$ & $+5.5$ & & \\
\bottomrule
\end{tabular}
\end{table}

Significant sex $\times$ genotype interactions were observed in L2/3 ($F(1,73) = 5.56$, $p < 0.05$) and L5 ($F(1,71) = 7.33$, $p < 0.01$). No effects of sex or genotype were found on UP state amplitude or frequency in any layer (all $F < 1.2$, all $p > 0.3$).

\subsection{mGluR5 Mediates Female-Specific Circuit Hyperexcitability}

Bath application of the mGluR5 negative allosteric modulator MTEP (3~$\mu$M, 2~hr) selectively rescued prolonged UP states in female Cre(+) slices (Table~\ref{tab:mtep}).

\begin{table}[H]
\centering
\caption{Effect of MTEP (3~$\mu$M) on UP state duration. Values are mean $\pm$ SEM. Statistical analysis by two-way ANOVA (drug $\times$ genotype) within each sex, with Sidak post-hoc tests. $n$ indicates number of slices from 4--8 mice per group.}
\label{tab:mtep}
\begin{tabular}{llcccccc}
\toprule
\textbf{Sex} & \textbf{Genotype} & \textbf{Vehicle (ms)} & \textbf{MTEP (ms)} & \textbf{$\Delta$ (ms)} & \textbf{$n$} & \textbf{$t$-stat} & \textbf{$p$-value} \\
\midrule
Female & Cre(+) & $648 \pm 41$ & $428 \pm 33$ & $-220$ & 38 & $4.12$ & $< 0.001$ \\
Female & Cre(-) & $418 \pm 29$ & $405 \pm 31$ & $-13$ & 34 & $0.31$ & $0.76$ \\
Male & Cre(+) & $425 \pm 32$ & $418 \pm 35$ & $-7$ & 36 & $0.14$ & $0.89$ \\
Male & Cre(-) & $402 \pm 30$ & $395 \pm 28$ & $-7$ & 32 & $0.17$ & $0.87$ \\
\bottomrule
\end{tabular}
\end{table}

Western blot analysis of L5-enriched cortical lysates revealed that mGluR5 protein levels were elevated selectively in female Cre(+) cortex (Table~\ref{tab:mglur5_protein}).

\begin{table}[H]
\centering
\caption{mGluR5 protein levels in L5-enriched cortical lysates, normalized to Cre(-) same-sex controls. Values are mean $\pm$ SEM of densitometry ratios. Statistical analysis by unpaired $t$-test within each sex.}
\label{tab:mglur5_protein}
\begin{tabular}{lcccccc}
\toprule
\textbf{Sex} & \textbf{Cre(-)} & \textbf{Cre(+)} & \textbf{Fold change} & \textbf{$n$ (mice)} & \textbf{$t$-stat} & \textbf{$p$-value} \\
\midrule
Female & $1.00 \pm 0.08$ & $1.42 \pm 0.11$ & $1.42$ & $6$ & $3.18$ & $< 0.01$ \\
Male & $1.00 \pm 0.09$ & $1.06 \pm 0.10$ & $1.06$ & $6$ & $0.45$ & $0.67$ \\
\bottomrule
\end{tabular}
\end{table}

\subsection{ER$\alpha$ Is Required for Female-Specific Hyperexcitability}

Two selective ER$\alpha$ antagonists rescued UP state duration in female Cre(+) mice. Genetic reduction of ER$\alpha$ via heterozygous Esr1 deletion corrected multiple phenotypes exclusively in females (Table~\ref{tab:eralpha}).

\begin{table}[H]
\centering
\caption{Genetic rescue by ER$\alpha$ reduction (Esr1$^{fl/+}$). Values normalized to same-sex Cre(-) controls (set to 1.0). Statistical analysis by one-way ANOVA with Tukey post-hoc tests within each sex. $n = 3$ mice per group, 17 sections per mouse.}
\label{tab:eralpha}
\begin{tabular}{llcccccc}
\toprule
\textbf{Measure} & \textbf{Sex} & \textbf{Cre(-)} & \textbf{Cre(+)} & \textbf{Cre(+)/Esr1$^{fl/+}$} & \textbf{$F$-stat} & \textbf{$p$ (Cre+ vs Cre-)} & \textbf{$p$ (rescue)} \\
\midrule
UP state (ms) & Female & $1.00 \pm 0.06$ & $1.54 \pm 0.09$ & $1.08 \pm 0.07$ & $14.2$ & $< 0.001$ & $< 0.01$ \\
UP state (ms) & Male & $1.00 \pm 0.07$ & $1.05 \pm 0.08$ & $1.03 \pm 0.07$ & $0.18$ & $0.84$ & $0.91$ \\
Protein synth. & Female & $1.00 \pm 0.05$ & $1.68 \pm 0.12$ & $1.12 \pm 0.08$ & $21.6$ & $< 0.001$ & $< 0.001$ \\
Protein synth. & Male & $1.00 \pm 0.06$ & $1.07 \pm 0.09$ & $1.04 \pm 0.07$ & $0.28$ & $0.72$ & $0.88$ \\
Soma size & Female & $1.00 \pm 0.04$ & $1.31 \pm 0.07$ & $1.06 \pm 0.05$ & $9.8$ & $< 0.01$ & $< 0.05$ \\
Soma size & Male & $1.00 \pm 0.05$ & $1.28 \pm 0.08$ & $1.25 \pm 0.07$ & $5.4$ & $< 0.05$ & $0.72$ \\
\bottomrule
\end{tabular}
\end{table}

Co-immunoprecipitation revealed elevated mGluR5--ER$\alpha$ complexes selectively in female cortex ($1.38 \pm 0.12$ vs. $1.00 \pm 0.09$ male, $t(10) = 2.54$, $p < 0.05$).

\subsection{ERK and Protein Synthesis Drive the Female Phenotype}

Pharmacological inhibition of ERK (U0126, 20~$\mu$M) and protein synthesis (anisomycin 20~$\mu$M; cycloheximide 60~$\mu$M) each rescued UP state duration selectively in female Cre(+) slices (Table~\ref{tab:pathway}).

\begin{table}[H]
\centering
\caption{Pharmacological rescue of female Cre(+) UP state duration by ERK and protein synthesis inhibitors. Values are mean $\pm$ SEM of L4 UP state duration. Statistical analysis by paired $t$-test (vehicle vs. drug within genotype). All recordings from female mice, $n = 12$--30 slices from 3--6 mice per condition.}
\label{tab:pathway}
\begin{tabular}{lcccccc}
\toprule
\textbf{Drug} & \textbf{Target} & \textbf{Cre(+) Vehicle} & \textbf{Cre(+) Drug} & \textbf{$\Delta$ (\%)} & \textbf{$t$-stat} & \textbf{$p$-value} \\
\midrule
U0126 (20 $\mu$M) & ERK/MEK & $638 \pm 42$ & $435 \pm 31$ & $-31.8$ & $3.74$ & $< 0.01$ \\
Anisomycin (20 $\mu$M) & Protein synth. & $651 \pm 39$ & $418 \pm 28$ & $-35.8$ & $4.62$ & $< 0.001$ \\
Cycloheximide (60 $\mu$M) & Protein synth. & $644 \pm 44$ & $422 \pm 33$ & $-34.5$ & $3.92$ & $< 0.01$ \\
MPP (1 $\mu$M) & ER$\alpha$ & $655 \pm 38$ & $441 \pm 30$ & $-32.7$ & $4.18$ & $< 0.001$ \\
GNE (1 $\mu$M) & ER$\alpha$ & $641 \pm 41$ & $448 \pm 35$ & $-30.1$ & $3.51$ & $< 0.01$ \\
\bottomrule
\end{tabular}
\end{table}

SUnSET assays revealed that de novo protein synthesis was elevated selectively in female Pten KO L5 neurons relative to neighboring PTEN-positive neurons ($1.68 \pm 0.12$ normalized fluorescence, $t(16) = 5.67$, $p < 0.0001$), with no change in male Cre(+) neurons ($1.07 \pm 0.09$, $t(16) = 0.78$, $p = 0.45$).

\subsection{Sensory Processing Deficits and Behavioral Abnormalities}

Female NSEPten KO mice showed deficits in temporal processing of auditory stimuli measured by gap-ASSR at P21 (Table~\ref{tab:behavior}).

\begin{table}[H]
\centering
\caption{Behavioral and sensory processing phenotypes in NSEPten KO mice. ITPC = inter-trial phase coherence during gap-ASSR. Social interaction measured in 3-chamber test. Open field assessed over 30~min. Fear conditioning measured as freezing percentage at 24~hr. Statistical analysis by two-way ANOVA (sex $\times$ genotype) with Sidak post-hoc tests.}
\label{tab:behavior}
\begin{tabular}{llccccc}
\toprule
\textbf{Measure} & \textbf{Sex} & \textbf{Cre(-)} & \textbf{Cre(+)} & \textbf{$F$(sex$\times$geno)} & \textbf{$p$(interaction)} & \textbf{$p$(post-hoc)} \\
\midrule
ITPC auditory cortex & Female & $0.42 \pm 0.04$ & $0.21 \pm 0.03$ & $6.24$ & $< 0.05$ & $< 0.01$ \\
ITPC auditory cortex & Male & $0.39 \pm 0.04$ & $0.28 \pm 0.04$ & --- & --- & $< 0.05$ \\
ITPC frontal cortex & Female & $0.38 \pm 0.05$ & $0.18 \pm 0.03$ & $5.81$ & $< 0.05$ & $< 0.01$ \\
ITPC frontal cortex & Male & $0.36 \pm 0.04$ & $0.26 \pm 0.04$ & --- & --- & $0.06$ \\
Social preference index & Female & $0.68 \pm 0.04$ & $0.44 \pm 0.05$ & $4.92$ & $< 0.05$ & $< 0.01$ \\
Social preference index & Male & $0.65 \pm 0.05$ & $0.58 \pm 0.05$ & --- & --- & $0.31$ \\
Open field distance (m) & Female & $42.1 \pm 3.2$ & $58.4 \pm 4.1$ & $5.18$ & $< 0.05$ & $< 0.01$ \\
Open field distance (m) & Male & $44.3 \pm 3.5$ & $47.8 \pm 3.8$ & --- & --- & $0.48$ \\
Freezing (\%) & Female & $38.2 \pm 4.1$ & $62.4 \pm 5.3$ & $7.14$ & $< 0.01$ & $< 0.001$ \\
Freezing (\%) & Male & $41.5 \pm 4.4$ & $44.8 \pm 4.6$ & --- & --- & $0.58$ \\
\bottomrule
\end{tabular}
\end{table}

Seizure susceptibility testing with flurothyl revealed mGluR5-dependent seizure severity selectively in females (Mantel--Cox survival analysis, $\chi^2 = 8.42$, $p < 0.01$ for female Cre(+) vs. Cre(-); $\chi^2 = 1.28$, $p = 0.26$ for males).

\subsection{Replication in the Germline Pten$^{+/-}$ Model}

The female-specific cortical hyperexcitability was replicated in the clinically relevant germline Pten$^{+/-}$ model at 6--8 weeks of age (Table~\ref{tab:germline}).

\begin{table}[H]
\centering
\caption{UP state recordings and seizure susceptibility in germline Pten$^{+/-}$ mice (6--8 weeks). UP state duration in L4 somatosensory cortex. Seizure survival assessed by repeated daily flurothyl at P90--100. Statistical analysis by two-way ANOVA (sex $\times$ genotype) and Mantel--Cox test.}
\label{tab:germline}
\begin{tabular}{llcccccc}
\toprule
\textbf{Measure} & \textbf{Sex} & \textbf{WT} & \textbf{Pten$^{+/-}$} & \textbf{$n$ (slices)} & \textbf{$n$ (mice)} & \textbf{$F$ or $\chi^2$} & \textbf{$p$-value} \\
\midrule
UP duration (ms) & Female & $425 \pm 22$ & $612 \pm 31$ & 58 & 12 & $F = 18.4$ & $< 0.001$ \\
UP duration (ms) & Male & $418 \pm 24$ & $441 \pm 28$ & 43 & 8 & $F = 0.42$ & $0.52$ \\
Total UP time (\%) & Female & $18.2 \pm 1.4$ & $28.6 \pm 2.1$ & 58 & 12 & $F = 14.7$ & $< 0.001$ \\
Total UP time (\%) & Male & $17.8 \pm 1.6$ & $19.4 \pm 1.8$ & 43 & 8 & $F = 0.48$ & $0.49$ \\
Seizure survival & Female & $92\%$ & $58\%$ & --- & 12 & $\chi^2 = 6.84$ & $< 0.01$ \\
Seizure survival & Male & $88\%$ & $79\%$ & --- & 10 & $\chi^2 = 0.82$ & $0.37$ \\
\bottomrule
\end{tabular}
\end{table}

\section{Discussion}

We have identified a sex-specific molecular mechanism by which deletion of the ASD-risk gene \textit{Pten} produces cortical circuit hyperexcitability selectively in females. The signaling cascade proceeds from Pten loss through mGluR5 upregulation, formation of mGluR5--ER$\alpha$ complexes, ERK activation, and enhanced de novo protein synthesis, ultimately prolonging spontaneous UP states and increasing seizure susceptibility. Several aspects of these findings merit discussion.

First, the female specificity of the circuit phenotype is striking given that Pten is deleted equally in both sexes. The critical sex-differential component is ER$\alpha$: genetic reduction of a single Esr1 allele rescues hyperexcitability, protein synthesis, and soma size exclusively in females. This indicates that ER$\alpha$ acts as a molecular switch that converts Pten loss into circuit hyperexcitability only when sufficient ER$\alpha$ signaling is present. Co-immunoprecipitation data suggest that physical interaction between mGluR5 and ER$\alpha$ is the proximal mechanism, consistent with prior reports of mGluR5--ER$\alpha$ coupling in non-neural tissues \citep{lopes2012interaction}.

Second, the involvement of mGluR5 connects this work to the broader landscape of ASD-related circuit dysfunction. Excessive mGluR5 signaling is the central pathophysiology in Fragile X syndrome \citep{huber2002altered, dolen2007correction}, and mGluR5-dependent UP state prolongation has been reported in Fmr1 KO mice \citep{hays2011altered, gibson2008imbalance}. Our finding that Pten deletion produces a qualitatively similar but sex-restricted mGluR5-dependent phenotype suggests convergent circuit-level pathophysiology across genetically distinct forms of ASD.

Third, the protein synthesis dependence of the female phenotype links circuit excitability to a well-established downstream effector of mGluR5 signaling \citep{huber2002altered, ronesi2012homer}. The SUnSET assay demonstrated cell-autonomous elevation of de novo protein synthesis in female Pten KO neurons \citep{schmidt2009sUnSET}, and this was corrected by ER$\alpha$ reduction, placing protein synthesis downstream of the mGluR5--ER$\alpha$ interaction.

Fourth, the replication of female-specific hyperexcitability and seizure vulnerability in germline Pten$^{+/-}$ mice is clinically significant. Germline PTEN mutations in humans produce variable phenotypes \citep{butler2005subset}, and our data suggest that sex may be an important modifier of disease severity.

The behavioral data reveal that female NSEPten KO mice exhibit deficits in temporal processing of auditory stimuli (gap-ASSR), reduced social interaction, hyperactivity, and enhanced fear learning, consistent with core and associated features of ASD \citep{lovelace2018mirroring}. The mGluR5 dependence of seizures in females but not males provides a mechanistic basis for the clinical observation that epilepsy is more prevalent in females with PTEN-related ASD.

A limitation of this study is that NSE-Cre targets predominantly excitatory neurons in layers III--V, so the contribution of inhibitory neuron Pten deletion cannot be assessed. Additionally, while we focused on somatosensory cortex, the phenotype may extend to other cortical regions.

\section{Methods}

\subsection{Animals}

All procedures were approved by the UT Southwestern IACUC. NSEPten KO mice were generated by crossing NSE-Cre mice with Pten$^{fl/fl}$ mice on a C57BL/6 background. For ER$\alpha$ reduction experiments, NSE-Cre/Pten$^{fl/fl}$ mice were crossed with Esr1$^{fl/+}$ mice. Germline Pten$^{+/-}$ mice were generated using CMV-Cre. Animals were maintained on Teklad Global 16\% Protein Rodent Diet to minimize phytoestrogens.

\subsection{Electrophysiology}

Thalamocortical slices (400~$\mu$m) were prepared from P18--25 (NSEPten KO) or 6--8 week old (Pten$^{+/-}$) mice. Slices were maintained in high-activity ACSF (5~mM KCl, 1~mM MgCl$_2$, 1~mM CaCl$_2$) at 32$^\circ$C. Extracellular multiunit recordings used 0.5~M$\Omega$ tungsten electrodes (FHC), amplified 10,000-fold, sampled at 2.5~kHz, and filtered at 500~Hz--3~kHz. UP states were detected using a 15$\times$ RMS noise threshold with a minimum duration of 200~ms and minimum inter-UP interval of 600~ms.

\subsection{Biochemistry}

Western blots used standard protocols with anti-mGluR5 primary antibody. Co-immunoprecipitation for mGluR5--ER$\alpha$ complexes used Protein A/G beads with anti-mGluR5 pulldown and anti-ER$\alpha$ detection. SUnSET assays for de novo protein synthesis followed \citet{schmidt2009sUnSET}. Soma size was measured from Nissl-stained sections using ImageJ.

\subsection{Behavioral Testing}

Gap-ASSR was recorded by EEG in P21 pups. Social interaction was assessed using the 3-chamber test. Open field locomotion was recorded for 30~min. Fear conditioning used a standard protocol with 24~hr retention test. Seizure susceptibility was assessed by flurothyl exposure with survival analysis (Mantel--Cox test).

\subsection{Statistical Analysis}

Two-way ANOVA (sex $\times$ genotype) was the primary statistical test. Post-hoc comparisons used Sidak correction. Unpaired or paired $t$-tests were used for two-group comparisons. Seizure data were analyzed by Mantel--Cox log-rank test. All data are presented as mean $\pm$ SEM. Significance threshold was $\alpha = 0.05$.

\section{Conclusion}

We have identified a female-specific mGluR5--ER$\alpha$--ERK--protein synthesis signaling pathway that drives cortical circuit hyperexcitability following \textit{Pten} deletion. This mechanism provides a molecular basis for sex-differential ASD phenotypes and identifies mGluR5 and ER$\alpha$ as potential sex-specific therapeutic targets for PTEN-related neurodevelopmental disorders.

\bibliographystyle{plainnat}
\bibliography{references}

\end{document}
